\begin{solapartat}
\begin{align*}
&m\omega^2(R_T+d) = \frac{GM_Tm}{(R_T+d)^2} \ \text{\punts{0,2}} \Rightarrow (R_T+d)^3 = \frac{GM_TT^2}{4\pi^2}\\
&\Rightarrow d = \sqrt[3]{\frac{GM_TT^2}{4\pi^2}} - R_T \ \text{\punts{0,6}} = 3,59\cdot 10^{4}\ \mathrm{km} \ \text{\punts{0,2}}
\end{align*}
Si deixen de restar el radi de la Terra se'ls resta \textbf{0,2} punts
\end{solapartat}

\begin{solapartat}
\[ E_c = \frac{1}{2}m\omega^2(R_T+d)^2 = \frac{1}{2}\frac{GM_Tm}{R_T+d} = 9,42\cdot 10^{9}\ \mathrm{J} \ \text{\punts{0,3}} \]
Per tal que el satèl·lit s'allunyi de l'atracció de la Terra, la seva energia mecànica ha de ser 0 \punts{0,2} $\Rightarrow$
\[ E_m = E_c + E_p = -\frac{1}{2}\frac{GM_Tm}{R_T+d} = -E_c \ \text{\punts{0,3}} \]
Per tant caldrà subministrar-li una energia igual a $E_c = 9,42\cdot 10^{9}\ \mathrm{J}$ \punts{0,2}
\end{solapartat}
